\subsection{The Shift to AI-Native Networking}

Nguyen et al.~\cite{nguyen20226g} outline the 6G vision of AI-native networks where DRL
manages massive, heterogeneous IoT connectivity in real time---conditions under which
hand-crafted optimisation policies are fundamentally insufficient.

\subsubsection{Deep Q-Networks in Communications}

The Deep Q-Network~\cite{mnih2015humanlevel} stabilises Q-learning with experience replay
and a periodically frozen target network. Experience replay breaks the temporal correlation
between consecutive UAV observations (important given the spatial correlation of the flight
path), while the target network prevents the Q-value feedback loop that causes divergence.
Luong et al.~\cite{luong2019drl} survey DRL in communications and observe that reward
shaping and state-space design are critical: poorly chosen representations prevent the agent
from learning the temporal patterns needed to anticipate future congestion. This directly
motivates the frame-stacking and urgency-proxy design choices in~\autoref{sec:method}.

\subsubsection{DRL for UAV Path Planning}

No prior work applies DRL specifically to UAV path planning in LoRaWAN networks with
explicit ADR and Capture Effect modelling. The closest precedents operate in the cellular
domain, where interference models are stationary and the MAC-layer dynamics (EMA-ADR,
duty cycles) that dominate LoRaWAN throughput are absent.

Challita et al.~\cite{challita2018interference} apply DRL to interference-aware path
planning for cellular-connected UAVs, showing that a learned policy can trade off wireless
latency against ground-user interference---achieving up to 33\% latency reduction over
shortest-path navigation and generalising to unseen network states. This trade-off is
directly analogous to the starvation--throughput balance addressed here, but their
stationary interference model does not capture the non-stationary EMA-ADR feedback
loop that makes the LoRaWAN setting harder.

Nemer et al.~\cite{nemer2022energyefficient} demonstrate that fairness-constrained DRL
(SBG-AC algorithm, Jain's Fairness Index as objective) outperforms energy-only DRL,
confirming that the reward function must explicitly penalise unequal service to prevent the
agent converging on policies that neglect underserved nodes. This finding directly motivates
the Age-of-Information penalty term in this work. However, Nemer et al.\ use fixed sensor
positions and a simplified channel model, leaving open whether the fairness gains transfer
to the stochastic, topology-varying LoRaWAN conditions that domain randomisation is
designed to address.

The absence of LoRaWAN-specific DRL prior work is not an oversight in this review: it
reflects a genuine gap in the literature, confirmed by the simulation tool analysis
(Section~\ref{sec:lit}), and constitutes the primary motivation for this
dissertation.

\subsubsection{Information Freshness and Domain Generalisation}

Kaul et al.~\cite{kaul2012realtime} show that the throughput-optimal update frequency is not
the freshness-optimal frequency, motivating the explicit AoI term in the reward function.
Abd-Elmagid et al.~\cite{abdelmagid2019aoi} demonstrate a fundamental AoI--throughput
trade-off in IoT networks, supporting the multi-term weighted reward design adopted here.

Tobin et al.~\cite{tobin2017domain} introduce Domain Randomisation to bridge the sim-to-real
gap, subsequently validated by OpenAI~\cite{openai2019dexterous} at scale in robotic
manipulation. Bengio et al.~\cite{bengio2009curriculum} show that curriculum learning---presenting
training examples by increasing difficulty---produces faster convergence and better
generalisation. Both techniques are combined in this work to produce a single policy
deployable across a 100$\times$ variation in operational area without retraining.
